% Глава: 16-virtual-memory
% Статус: наполнено (первый прогон).
\chapter{Виртуальная память}
\label{ch:16-virtual-memory}

\section{Зачем нужна виртуальная память}

Каждому процессу хотелось бы видеть «свою» память, начиная
с адреса 0, и не сталкиваться с другими процессами.
\textbf{Виртуальная память} --- это механизм ОС, который
даёт каждому процессу \emph{отдельное адресное пространство},
а ядро переводит виртуальные адреса в физические.

Это же самое происходит на вашей ВМ: каждый процесс,
запущенный на \fpath{pglab.nsu.ru}, живёт в собственной
виртуальной памяти.

Выгоды:
\begin{itemize}
  \item \textbf{Изоляция} --- сбой в одном процессе не
    «ломает» соседний;
  \item \textbf{Больше памяти, чем в RAM} --- неиспользуемые
    страницы выгружаются на диск (\emph{swap});
  \item \textbf{Общий код} --- один и тот же текст библиотеки
    может жить в RAM один раз и «видеться» многим процессам.
\end{itemize}

\section{Строение адресного пространства процесса}

У процесса есть несколько областей:

\begin{figure}[htb]
  \centering
  \begin{tikzpicture}[node distance=0.6cm]
    \node[box, minimum width=7cm, minimum height=1cm] (stack) {stack (растёт вниз)};
    \node[box, minimum width=7cm, minimum height=1cm, below=of stack] (heap) {heap (malloc, растёт вверх)};
    \node[box, minimum width=7cm, minimum height=1.6cm, below=of heap] (data) {data / bss\\ (глобальные переменные)};
    \node[box, minimum width=7cm, minimum height=1cm, below=of data] (text) {text (код программы)};
  \end{tikzpicture}
  \caption{Основные области адресного пространства процесса.}
  \label{fig:mem-layout}
\end{figure}

\begin{itemize}
  \item \textbf{stack} --- локальные переменные и адреса
    возврата; управляется автоматически, ограничен
    (обычно несколько мегабайт);
  \item \textbf{heap} --- память, запрашиваемая через
    \cmd{malloc}/\cmd{new}; управляется программой;
  \item \textbf{data/bss} --- глобальные и статические
    переменные;
  \item \textbf{text} --- код программы (только чтение).
\end{itemize}

\section{Страницы и таблица страниц}

Память разбита на \textbf{страницы} (обычно 4 КБ).
Ядро хранит для процесса \textbf{таблицу страниц}:
соответствие «виртуальная страница $\to$ физическая
страница (или файл на диске)». При обращении к адресу
аппарат (\emph{MMU}/TLB) через таблицу переводит адрес;
если страницы нет в RAM --- происходит \textbf{page fault},
ядро подгружает её.

\begin{figure}[htb]
  \centering
  \begin{tikzpicture}[node distance=1.4cm]
    \node[process] (cpu) {CPU\\[0.2em] виртуальный адрес};
    \node[process, below=of cpu] (mmu) {MMU / TLB\\[0.2em] перевод адреса};
    \node[process, below=of mmu] (pt) {таблица страниц};
    \node[storage, below left=1.3cm and -0.6cm of pt] (ram) {RAM};
    \node[storage, below right=1.3cm and -0.6cm of pt] (disk) {диск (swap)};
    \draw[flow, <->] (cpu) -- (mmu);
    \draw[flow, <->] (mmu) -- (pt);
    \draw[flow, <->] (pt) -- (ram);
    \draw[flow, <->] (pt) -- (disk);
  \end{tikzpicture}
  \caption{Перевод виртуальных адресов в физические.}
  \label{fig:vm}
\end{figure}

\section{\cmd{mmap} и общая память}

Функция \cmd{mmap} позволяет «пришить» файл или
анонимную область к адресному пространству процесса.
Через неё реализованы и загрузка библиотек, и
\textbf{общая память} между процессами (см.\
\nameref{ch:18-ipc}): два процесса могут видеть одну
и ту же физическую страницу.

\section{Связь с процессами и Docker}

Каждый процесс (\nameref{ch:12-processes}) имеет своё
адресное пространство; когда процесс вызывается из
оболочки, у него отдельный стек и куча. Контейнеры
Docker (\nameref{ch:19-docker-basics}) разделяют
\emph{имена} и ресурсы, но виртуальная память процессов
внутри контейнера работает по тем же правилам;
ограничения по RAM на контейнер --- через cgroups
(\nameref{ch:23-docker-resources}).

\section{Типичные проблемы, связанные с памятью}

\begin{itemize}
  \item \textbf{Утечка} --- \cmd{malloc} без \cmd{free};
    память накапливается, процесс «раздувается»;
  \item \textbf{Переполнение стека} --- глубокая
    рекурсия или огромный локальный массив;
  \item \textbf{Use-after-free} --- обращение к памяти
    после \cmd{free};
  \item \textbf{Дangling pointer} --- указатель на
    объект, который уже уничтожен.
\end{itemize}

Всё это ищут инструментами из \nameref{ch:17-debugging}.

\section{Задачи для самопроверки}
\begin{enumerate}
  \item Зачем нужна виртуальная память? Назовите
    три выгоды.
  \item Перечислите основные области адресного
    пространства и что в них хранится.
  \item Что такое страница и таблица страниц?
  \item Что происходит при page fault?
  \item Чем heap отличается от stack по управлению
    и размеру?
\end{enumerate}
